\section{附录：术语与缩写}
\subsection{主要术语解读}
\begin{table}[H]
\centering
\begin{tabular}{p{4cm}p{11cm}}
\toprule
术语 & 解释 \\
\midrule
Chain-of-Thought & 让模型列出中间推理步骤，提供 self-check 的基础。 \\
Self-Rewarding & 模型自行生成任务并利用 reward 来提升 reasoning，不依赖人工 label。 \\
Self-Instruct & 从少量 seed 自动生成 instruction→input→output，扩充训练语料。 \\
Verified Response & 模型自我检查后生成的最终回答，作为 reward 的可靠信号。 \\
Meta Judge & 比较 draft 与 verified，决定哪个更可信的 evaluator。 \\
SelfT Evaluators & 以 verified response 为标签训练的 judge 数据集。 \\
\bottomrule
\end{tabular}
\caption{本讲笔记里反复提及的几个术语。}
\end{table}

\subsection{相关缩写}
\begin{itemize}
    \item PRM（Process Reward Model）：逐步奖励，用于鼓励每一步 reasoning。
    \item ORM（Outcome Reward Model）：只看最终答案，Meta judge 典型做法。
    \item RLHF（Reinforcement Learning from Human Feedback）：传统含人工反馈的 reward 训练。
    \item DPO（Direct Preference Optimization）：一种无 RL，也有 preference learning 的参数更新策略。
\end{itemize}

\subsection{本章小结}
术语与缩写的附录帮助复盘讲座中的关键词，便于快速定位 Self-Rewarding、Meta Judge、Self-Instruct 之间的关系。

